%! Inverses of Exponential Functions — Unit 5, Lesson 5.2 — Lesson Plan
\documentclass[10pt]{article}
\usepackage{precalculus-article}
\usepackage{precalculus-boxes}
\usepackage{graphicx}
\graphicspath{{images/}}
\newcommand{\TallMath}[1]{$\displaystyle #1\rule[-1.4em]{0pt}{3.2em}$}
\newcommand{\CourseName}{Precalculus}
\newcommand{\SchoolYear}{2026--2027}
\newcommand{\MeetingLength}{55 minutes}
\newcommand{\UnitNumberName}{Unit 5: Logarithmic Functions \quad}
\newcommand{\LessonNumberName}{Lesson 5.2: Inverses of Exponential Functions}

\begin{document}
\begin{center}
    {\Large\bfseries \CourseName: \SchoolYear \\
    {\small\bfseries \UnitNumberName \LessonNumberName}}
\end{center}

\begin{tcolorbox}[colback=sky, colframe=navy, boxrule=0.9pt, arc=2mm,
    left=2mm, right=2mm, top=1.5mm, bottom=1.5mm]
    \textbf{Primary Objective:} Students construct the logarithmic function as the inverse of
    an exponential function (initial value 1), represent the inverse using tables and algebraic
    notation, verify the inverse relationship through composition, and connect the
    additive/multiplicative growth patterns of each function type.
    \hfill\textit{\small Skills: AP 1.C, AP 2.B \quad LO 2.10.A}
\end{tcolorbox}

\renewcommand{\arraystretch}{1.6}
\begin{skillbox}[Priority Ideas \& Skills]{goldbox}
\begin{tabularx}{\linewidth}{@{}X X@{}}
\textbf{AP Skills} &
\textbf{Key Understandings} \\
\midrule
\textbf{AP 1.C} --- Construct new functions using inverses. &
\begin{itemize}[leftmargin=1.2em,itemsep=2pt,topsep=0pt]
  \item $f(x) = \log_b x$ and $g(x) = b^x$ are inverse functions:
        $g(f(x)) = f(g(x)) = x$. (EK 2.10.A.3)
  \item If $(s,t)$ is on $g(x) = b^x$, then $(t,s)$ is on $f(x) = \log_b x$.
        (EK 2.10.A.5)
\end{itemize} \\
\textbf{AP 2.B} --- Construct equivalent representations. &
\begin{itemize}[leftmargin=1.2em,itemsep=2pt,topsep=0pt]
  \item Exponential output changes multiplicatively as input changes additively;
        logarithmic output changes additively as input changes multiplicatively.
        (EK 2.10.A.2)
  \item General form: $f(x) = a\log_b x$, $b > 0$, $b \ne 1$, $a \ne 0$.
        (EK 2.10.A.1)
\end{itemize} \\
\end{tabularx}
\end{skillbox}

\begin{skillbox}[{Vocabulary, Concepts, \& Theorems}]{greenbox}
\renewcommand{\arraystretch}{1.4}
\begin{tabularx}{\linewidth}{@{} >{\leavevmode\bfseries\color{navy}}p{0.28\linewidth} X @{}}
Logarithmic function & A function of the form $f(x) = a\log_b x$ where $b > 0$, $b \ne 1$,
                       and $a \ne 0$; the inverse of a general-form exponential function with
                       initial value 1. \\
Base of a logarithm  & The value $b$ in $\log_b x$; the same base as the corresponding
                       exponential $b^x$. \\
Inverse functions    & Two functions $f$ and $g$ such that $f(g(x)) = g(f(x)) = x$ for all
                       $x$ in the appropriate domain. \\
Additive change      & Output values that increase (or decrease) by a constant amount across
                       equal intervals; characteristic of logarithmic functions. \\
Multiplicative change & Output values that increase by a constant factor across equal intervals;
                       characteristic of exponential functions. \\
Reflection over $y=x$ & The geometric relationship between a function and its inverse: the
                       graph of $f(x) = \log_b x$ is the reflection of $g(x) = b^x$ over the
                       line $y = x$. (EK 2.10.A.4) \\
\end{tabularx}
\end{skillbox}

\begin{skillbox}[Activate Prior Knowledge \& Spiral Review]{sky}
    Spiral review (warm-up) covers: evaluating $\log_2 8$ and $2^3$ to notice they share the
    same result; completing $\log_3(3^5)$ and $3^{\log_3 7}$ to surface the cancellation
    property; and identifying the reflected point of $(3, 8)$ on an inverse function.
    These connect directly to EK 2.10.A.3 (composition) and EK 2.10.A.5 (swapped pairs).
\end{skillbox}

\begin{skillbox}[Hook]{sky}
    Write on the board: \textbf{A population model is $P(t) = 2^t$.}

    Ask: ``How long until the population reaches 64?'' Have students guess, then try to solve
    algebraically. Most will notice $2^t = 64 \Rightarrow t = 6$, but ask: \textbf{``Is there
    a function that gives us $t$ directly, given $P$?''} Pose the question: ``If exponentiation
    is the operation, what is the inverse operation?'' This motivates the logarithm as the
    function that ``undoes'' an exponential.
\end{skillbox}

\begin{skillbox}[Lesson]{sky}
\begin{multicols}{2}
\textbf{Part 1 (12 min): Growth Patterns Side-by-Side.}

Build a table for $g(x) = 2^x$ at $x = 0,1,2,3$. Ask: as $x$ increases by 1, what happens
to $g(x)$? (Multiplies by 2.) Now build the reflected table for $f(x) = \log_2 x$ using the
swapped pairs $(1,0),(2,1),(4,2),(8,3)$. Ask: as $x$ multiplies by 2, what happens to $f(x)$?
(Adds 1.) Have students fill the guided-notes table and write the additive/multiplicative
sentences.

\columnbreak

\textbf{Part 2 (12 min): Inverse Relationship.}

Define: $f$ and $g$ are inverses when $f(g(x)) = g(f(x)) = x$. Verify:
$g(f(x)) = 2^{\log_2 x} = x$ and $f(g(x)) = \log_2(2^x) = x$.

Emphasize EK 2.10.A.5: if $(s,t)$ is on $g(x) = b^x$, then $(t,s)$ is on $f(x) = \log_b x$.
Work through the $g(x) = 3^x$ inverse-pair exercise in the notes together.

\textbf{Part 3 (8 min): General Form.}

Introduce $f(x) = a\log_b x$ as the general form (EK 2.10.A.1). Show that the vertical
dilation by $a$ still preserves the inverse relationship with $g(x) = b^x$ when $a = 1$, but
stretches/compresses the outputs otherwise. Return to the hook: $t = \log_2 P$, so $t(64) = 6$.
\end{multicols}
\end{skillbox}

\begin{skillbox}[Active Monitoring]{redbox}
\begin{itemize}[leftmargin=1.2em,itemsep=2pt,topsep=0pt]
  \item After Part 1: cold-call ``How are the input/output patterns of $g(x)=2^x$ and
        $f(x)=\log_2 x$ opposite?'' Expect: exponential multiplies output; log adds to output.
  \item Watch for students who write $\log_b(b^x) = b$ instead of $x$ --- address immediately
        by tracing through the composition step by step.
  \item After Part 2: ask ``If $(5, 32)$ is on $g(x) = 2^x$, what point is on $f(x) = \log_2 x$?''
        Expect $(32, 5)$. Students who say $(2,5)$ are confusing the base with the input.
  \item During Part 3: ensure students see $a\log_b x$ as a vertical dilation and not a change
        of base or a new kind of function.
\end{itemize}
\end{skillbox}

\begin{skillbox}[Group Work \& Differentiation]{redbox}
\begin{itemize}[leftmargin=1.2em,itemsep=2pt,topsep=0pt]
  \item \textbf{Tier R (Remediate):} Complete a table for $g(x) = 4^x$; swap pairs to form the
        inverse table; identify the inverse as $\log_4 x$; evaluate from the table; verify
        composition once.
  \item \textbf{Tier A (Apply):} Inverse of $g(t) = 5^t$ in the bacteria context; composition
        verification; evaluate in context (time to reach populations of 125 and 625).
  \item \textbf{Tier E (Extension):} Compare $f(x) = \log_2 x$ and $h(x) = 3\log_2 x$; build
        tables; check additive-growth property for both; determine when $f(x) = h(x)$.
\end{itemize}
\end{skillbox}

\begin{skillbox}[Individual Work \& Assessment]{redbox}
Exit ticket (3 questions, 1 page):
\begin{enumerate}[leftmargin=1.5em,itemsep=2pt,topsep=2pt]
  \item If $(2,9)$ is on $g(x)=b^x$, what point is on $f(x)=\log_b x$? (Tests EK 2.10.A.5.)
  \item Compute $g(f(27))$ where $f(x) = \log_3 x$, $g(x) = 3^x$. (Tests EK 2.10.A.3.)
  \item Complete a short table for $f(x) = \log_2 x$ using the inverse relationship.
        (Tests EK 2.10.A.2.)
\end{enumerate}

Sort by result: students missing \#1 need work on swapped pairs (EK 2.10.A.5);
missing \#2 need composition practice; missing \#3 need the additive growth pattern
(EK 2.10.A.2).
\end{skillbox}

\begin{skillbox}[Reinforcement \& Extension]{goldbox}
\textbf{Homework (6 problems):} write the inverse of $g(x) = 6^x$ and complete a paired table;
verify $f(x) = \log_5 x$ and $g(x) = 5^x$ are inverses using both compositions; contextual
inverse with $A(t) = 10^t$ (algae area); written explanation of additive vs.\ multiplicative
growth patterns; evaluate $h(x) = 2\log_3 x$; AP-style MC on graph reflection.

\textbf{Preview of Lesson 2.11:} Logarithmic functions as function families --- domain, range,
concavity, end behavior (as $x \to 0^+$ and $x \to \infty$), and why $f(x) = \log_b x$ is
always strictly increasing (for $b > 1$) or strictly decreasing (for $0 < b < 1$).
\end{skillbox}

\end{document}
